\documentclass[11pt,letterpaper]{article}
\usepackage[utf8]{inputenc}
\usepackage[T1]{fontenc}
\usepackage[spanish,es-tabla]{babel}
\usepackage{lmodern}
\usepackage[margin=2.5cm]{geometry}
\usepackage{amsmath,amssymb,bm}
\usepackage{booktabs,array}
\usepackage{graphicx}
\usepackage{float}
\usepackage{enumitem}
\usepackage{tikz}
\usetikzlibrary{backgrounds,fit,positioning}
\usepackage{quantikz}
\usepackage[colorlinks=true,linkcolor=blue!50!black,citecolor=blue!50!black,urlcolor=blue!50!black]{hyperref}
\usepackage{caption}
\captionsetup{font=small,labelfont=bf}

\providecommand{\ket}[1]{\left|#1\right\rangle}
\providecommand{\bra}[1]{\left\langle#1\right|}
\providecommand{\braket}[2]{\left\langle#1\middle|#2\right\rangle}
\newcommand{\R}{\mathbb{R}}

\title{\textbf{Propuesta de proyecto final:}\\[2pt]
Diferenciación de perros y gatos con aprendizaje automático cuántico}
\author{Computación Cuántica -- Proyecto Final\\
Prof.\ Naomi Itzel Reyes Granados \;\;|\;\; Ayud.\ Sebastián González Juárez}
\date{Entrega: 7 de diciembre de 2026}

\begin{document}
\maketitle

\begin{abstract}
\noindent Esta propuesta plantea clasificar imágenes de perros y gatos con un clasificador
variacional cuántico (VQC) y compararlo contra modelos clásicos. El conjunto de datos es
\emph{Dogs vs.\ Cats} de Kaggle. Como las imágenes no pueden entrar directamente a un circuito
de pocos qubits, se usa un extractor de características clásico preentrenado que reduce cada imagen
a un vector pequeño, el cual después se codifica en el circuito. La propuesta sigue la metodología de
\texttt{ProyectosCC.pdf}: preprocesamiento sin fuga de información, reducción con PCA, comparación
contra clasificadores clásicos y métricas de accuracy, matriz de confusión y F1. Adopta el enfoque de
\emph{data re-uploading} sugerido en la Opción 1 (Clasificación de Semillas).
\end{abstract}

\tableofcontents
\bigskip

%=====================================================================
\section{Narrativa y problema real}
%=====================================================================
Un refugio de animales recibe a diario cientos de fotografías de animales encontrados o dados en
adopción, muchas de ellas tomadas con teléfonos móviles y sin etiquetar. Clasificar manualmente
cada imagen como perro o gato es lento y consume tiempo del personal veterinario y de adopción. Un
sistema automático de primera clasificación permite ordenar el catálogo, acelerar el registro y
detectar errores de carga antes de publicar un animal en la plataforma de adopción.

El problema es binario: dada una imagen $\bm{I}$, decidir si corresponde a un perro ($y=1$) o a un
gato ($y=0$). Aunque la tarea parece sencilla para una persona, es un reto para modelos pequeños
porque las fotografías varían en iluminación, fondo, pose, encuadre y tamaño del animal.

\subsection{Pregunta de investigación y objetivos}
\begin{quote}
\emph{¿Puede un clasificador variacional cuántico de pocos qubits, alimentado con
características extraídas por una red neuronal preentrenada, alcanzar un desempeño comparable a
clasificadores clásicos de capacidad similar en la distinción perro--gato?}
\end{quote}
Objetivos específicos: (i) formalizar el problema y definir el preprocesamiento sin fuga de
información; (ii) extraer y reducir las características de las imágenes a un tamaño compatible con
la codificación cuántica; (iii) implementar y entrenar un VQC con \emph{data re-uploading};
(iv) compararlo con baselines clásicos mediante accuracy, F1, matriz de confusión y curva ROC/AUC;
(v) analizar la influencia de la profundidad del circuito y del número de \emph{shots}; (vi)
discutir honestamente si existe o no ventaja cuántica.

\subsection{Formalización}
Sea $\mathcal{D}=\{(\bm{I}_i,y_i)\}_{i=1}^{N}$ con $y_i\in\{0,1\}$. Se define un extractor
$g:\bm{I}\mapsto\R^{2048}$ (preentrenado y congelado) y una proyección $\pi:\R^{2048}\to\R^{d}$ con
$d=4$ (número de qubits). El clasificador es
$f_{\bm\theta}:\R^{d}\to[0,1]$, con $f_{\bm\theta}(\bm{x})=P(y=1\mid\bm{x})$ y $\bm{x}=\pi(g(\bm{I}))$,
que minimiza la entropía cruzada binaria
\begin{equation}
\mathcal{L}(\bm\theta)=-\frac{1}{N}\sum_{i=1}^{N}\Big[y_i\log f_{\bm\theta}(\bm{x}_i)+(1-y_i)\log\big(1-f_{\bm\theta}(\bm{x}_i)\big)\Big].
\label{eq:loss}
\end{equation}
Como $\bm\theta$ contiene ángulos de rotación, el problema es \textbf{no convexo}; por ello se
reportan resultados con varias semillas de inicialización.

%=====================================================================
\section{Datos}
%=====================================================================
Se usa el conjunto \emph{Dogs vs.\ Cats} de la competencia de Kaggle (\cite{kaggle2013}): 25\,000
imágenes a color de tamaño variable, etiquetadas como perro o gato, con distribución balanceada
(12\,500 por clase). Las imágenes de prueba de la competencia no tienen etiqueta pública, por lo
que la evaluación se hace con una partición del conjunto etiquetado.

\paragraph{Partición.} Partición estratificada 70\,\% entrenamiento / 15\,\% validación /
15\,\% prueba, con \texttt{random\_state=42}. El conjunto de prueba se usa una sola vez al final.
Con 25\,000 imágenes se obtienen aproximadamente 17\,500 / 3\,750 / 3\,750 ejemplos.

\paragraph{Preprocesamiento.} Todas las imágenes se redimensionan a $224\times224$ y se normalizan
con la media y desviación estándar de ImageNet. Esta transformación no depende de los datos, por lo
que no introduce fuga de información. Las etapas que sí dependen de los datos (PCA y escalado a
$[0,\pi]$) se ajustan únicamente con el conjunto de entrenamiento.

\paragraph{Balance.} Al ser balanceado, la accuracy es una métrica confiable; aun así se reporta F1
para la clase positiva y la matriz de confusión.

%=====================================================================
\section{Pipeline de extracción de características}
%=====================================================================
Codificar una imagen completa con un qubit por píxel es imposible en hardware actual. Por ello se
adopta un enfoque híbrido clásico-cuántico:
\begin{enumerate}[leftmargin=*]
\item \textbf{Extractor congelado.} Una red convolucional preentrenada en ImageNet (por ejemplo,
ResNet-18) se usa sin reentrenar y produce un vector de 512 dimensiones por imagen (capa de pooling
global). Esto aprovecha el conocimiento visual general sin costo de entrenamiento.
\item \textbf{Estandarización.} Media 0 y varianza 1, ajustada con entrenamiento.
\item \textbf{PCA a 4 componentes.} Reduce 512 a 4 dimensiones, una por qubit. Se ajusta con
entrenamiento y se aplica a validación y prueba. La varianza retenida se reporta como criterio de
selección de $d$ (se evaluará también $d\in\{4,6,8\}$ como ablación de costo).
\item \textbf{Escalado a $[0,\pi]$} con min--max ajustado en entrenamiento. Como $R_X(\alpha x)$ es
$2\pi$-periódica en el ángulo, este rango mantiene una correspondencia inyectiva entre valores y
estados en la esfera de Bloch.
\end{enumerate}
Con esta estructura, el circuito es el único componente cuántico y el resto es clásico y
reproducible.

%=====================================================================
\section{Enfoque cuántico: clasificador de \emph{data re-uploading}}
%=====================================================================
\subsection{Fundamento}
Un modelo variacional evalúa $f(\bm{x})=\bra{0}U^\dagger(\bm{x},\bm\theta)\,M\,U(\bm{x},\bm\theta)\ket{0}$.
Si el dato se codifica una sola vez con puertas $e^{-ixH/2}$, el modelo es una serie de Fourier con
frecuencias $\{-1,0,1\}$ \cite{schuld2021}, demasiado poco expresiva para fronteras complejas. La
técnica de \emph{data re-uploading} \cite{perez2020} intercala $L$ veces codificación y bloques
entrenables:
\begin{equation}
U_{\bm\theta}(\bm{x})=\prod_{l=L}^{1}\Big[W_l(\bm\theta_l)\,S_l(\bm{x};\bm\alpha_l)\Big],\qquad
S_l(\bm{x};\bm\alpha_l)=\bigotimes_{q=1}^{4}R_X(\alpha_{l,q}\,x_q),
\end{equation}
donde $W_l$ combina una rotación genérica $\mathrm{Rot}(\phi,\vartheta,\omega)$ por qubit y un anillo
de CNOT, y los $\alpha_{l,q}$ son escalas entrenables. Así, cada re-subida amplía el espectro de
frecuencias accesible y permite fronteras de decisión más complejas.

\subsection{Lectura binaria}
Para clasificación binaria se define la salida como la probabilidad de un subconjunto de estados
base. Por ejemplo, se asigna la clase perro al conjunto de estados con el qubit de salida en $\ket{1}$:
\begin{equation}
p_1(\bm{x})=\sum_{s:\,s_0=1}\big|\braket{s}{\psi(\bm{x})}\big|^{2},\qquad
f_{\bm\theta}(\bm{x})=\sigma\big(\beta\,(2p_1(\bm{x})-1)\big),
\end{equation}
con $\sigma$ la función logística y $\beta>0$ entrenable, que evita que la salida quede casi
constante (como en la propuesta de semillas, donde la temperatura se volvió necesaria).

\subsection{Número de parámetros}
Con $n=4$ qubits y $L$ capas, $|\bm\theta|=4nL+1$ (rotaciones, escalas y $\beta$). Para $L=4$ se
tienen 65 parámetros, cantidad comparable a la de un perceptrón multicapa pequeño.

\subsection{Optimización y costo}
Se minimiza \eqref{eq:loss} con Adam (tasa inicial $0.05$, decaimiento coseno, \emph{minibatch} de
128, 20 épocas) sobre simulador de vector de estado, con gradientes por diferenciación automática.
En hardware real se usaría \emph{parameter-shift}, que requiere dos evaluaciones del circuito por
parámetro y por muestra; para $L=4$ y $S=1000$ shots, el entrenamiento en hardware costaría del
orden de $10^{10}$ shots, lo que se reporta como análisis de costo.

\subsection{¿Por qué podría resolver el problema?}
\begin{enumerate}[leftmargin=*]
\item \textbf{Expresividad.} Las re-subidas permiten fronteras no lineales en el espacio de
características.
\item \textbf{Entrelazamiento.} Los CNOT generan dependencias entre productos de funciones
trigonométricas de distintas componentes, análogas a capas ocultas.
\item \textbf{Universalidad.} Con suficientes capas la familia puede aproximar funciones continuas
acotadas \cite{perez2020}; sin embargo, la universalidad no garantiza que un optimizador local la
encuentre, por lo que la evidencia empírica es parte esencial de la justificación.
\end{enumerate}

%=====================================================================
\section{Librerías y justificación de su uso}
%=====================================================================
Según el enunciado, las librerías extra deben justificarse:
\begin{itemize}[leftmargin=*]
\item \textbf{PyTorch y torchvision:} cargan el extractor ResNet-18 preentrenado y transforman las
imágenes. No existe una alternativa razonable para obtener las características sin reimplementar
una red completa.
\item \textbf{PennyLane:} define el circuito variacional y calcula gradientes automáticamente.
\item \textbf{Qiskit:} verificación cruzada del circuito y estimación de profundidad y compuertas.
\item \textbf{scikit-learn:} partición estratificada, PCA, estandarización, baselines clásicos y
métricas.
\end{itemize}

%=====================================================================
\section{Baselines clásicos y diseño experimental}
%=====================================================================
\paragraph{Métricas.} \emph{Accuracy}; F1 (clase perro); matriz de confusión; área bajo la curva
ROC (AUC). Todas se reportan con media y desviación estándar sobre tres semillas.

\paragraph{Baselines} (mismas particiones y mismas 4 componentes PCA, salvo que se indique):
\begin{itemize}[leftmargin=*]
\item Regresión logística.
\item SVM con kernel RBF.
\item MLP pequeño de 6 neuronas (capacidad comparable al VQC).
\item Random Forest.
\item Con las 512 características originales: regresión logística y SVM-RBF, para medir cuánto se
pierde al comprimir a 4 qubits.
\end{itemize}

\paragraph{Ablaciones.} (a) Profundidad $L\in\{1,2,4,6\}$; (b) VQC sin re-subida de datos
($L=4$, datos codificados solo en la primera capa); (c) número de componentes $d\in\{4,6,8\}$ con
el costo de qubits asociado.

\paragraph{Muestreo finito.} Evaluación del modelo entrenado con $S\in\{50,200,1000,5000\}$ shots
sobre 1\,000 imágenes de prueba, comparada con la evaluación exacta.

%=====================================================================
\section{Resultados esperados y criterio de éxito}
%=====================================================================
Se espera que los clasificadores clásicos sobre 512 características superen el 95\,\% de accuracy,
dado que el extractor ya separa bien ambas clases. Se considera exitoso el VQC si:
\begin{itemize}[leftmargin=*]
\item alcanza una accuracy dentro de 3 puntos porcentuales del mejor clásico con la misma entrada
de 4 componentes;
\item supera claramente al clasificador de centroides o a un modelo lineal sin no linealidad, lo
que indicaría que el circuito aporta capacidad no lineal;
\item muestra degradación controlada al usar pocos shots, con análisis del costo.
\end{itemize}
Un resultado negativo (VQC por debajo de los clásicos) también es un resultado válido y se
discutiría como parte de las conclusiones.

%=====================================================================
\section{Conclusiones esperadas}
%=====================================================================
\begin{enumerate}[leftmargin=*]
\item \textbf{Cuello de botella.} La compresión de 512 a 4 características es el principal limitante;
se espera que el desempeño dependa más de la entrada que del clasificador cuántico.
\item \textbf{Ventaja cuántica.} Es poco probable que exista ventaja: un espacio de Hilbert de dimensión
16 se simula trivialmente. El aporte es caracterizar un VQC compacto en una tarea visual real y sus costos.
\item \textbf{Costo práctico.} El VQC requiere mucho más tiempo de simulación y, en hardware, un
número de mediciones que crece con el tamaño del conjunto de entrenamiento.
\end{enumerate}

%=====================================================================
\section{Cronograma ajustado a la pre-entrega y la entrega}
%=====================================================================
El proyecto se organiza alrededor de dos hitos: la \textbf{pre-entrega} (presentación en PDF de hasta
20 minutos, 23 de noviembre de 2026) y la \textbf{entrega final} (Jupyter Notebook ejecutable y
documento, 7 de diciembre de 2026). Antes de la pre-entrega debe validarse el proyecto con la
profesora, porque no pertenece a las cuatro propuestas del enunciado.

\begin{table}[H]
\centering\small
\begin{tabular}{@{}p{3.3cm}p{8.9cm}p{2.3cm}@{}}
\toprule
Fechas & Actividad & Hito \\ \midrule
14--18 nov & Validar el proyecto con la profesora; investigación teórica; formalización y función
objetivo; EDA del dataset. & Validación \\
19--22 nov & Extracción de características con el extractor congelado; baselines clásicos; primeras
pruebas del circuito; preparación de diapositivas. & Avances de código \\
23 nov & Presentación de pre-entrega (todos los integrantes participan). & \textbf{Pre-entrega} \\
24--29 nov & Implementación del VQC en PennyLane y verificación con Qiskit; entrenamiento con
tres semillas. & Modelo entrenado \\
30 nov--3 dic & Ablaciones (profundidad, re-subida, $d$); muestreo finito; análisis de costo. & Resultados \\
4--6 dic & Cuaderno Jupyter final reproducible; redacción de conclusiones; revisión contra la rúbrica. & Cierre \\
7 dic & Entrega final del cuaderno y del documento. & \textbf{Entrega} \\
\bottomrule
\end{tabular}
\end{table}

%=====================================================================
\section{Riesgos y mitigación}
%=====================================================================
\begin{itemize}[leftmargin=*]
\item \textbf{Tiempo de simulación:} usar $d=4$ y un subconjunto de entrenamiento si es necesario,
documentando el cambio.
\item \textbf{Mesetas estériles:} inicializaciones múltiples y arquitectura poco profunda.
\item \textbf{Extractor sesgado:} el modelo ImageNet puede tener sesgos; se reporta el análisis de
errores por tipo de imagen.
\end{itemize}

%=====================================================================
\begin{thebibliography}{9}
\bibitem{kaggle2013} Kaggle. \emph{Dogs vs.\ Cats}. Competencia de clasificación de imágenes,
2013. \url{https://www.kaggle.com/c/dogs-vs-cats}.
\bibitem{perez2020} A.\ Pérez-Salinas, A.\ Cervera-Lierta, E.\ Gil-Fuster y J.\ I.\ Latorre.
\emph{Data re-uploading for a universal quantum classifier}. Quantum, 4:226, 2020.
\bibitem{schuld2021} M.\ Schuld, R.\ Sweke y J.\ J.\ Meyer. \emph{Effect of data encoding on the
expressive power of variational quantum-machine-learning models}. Phys.\ Rev.\ A, 103:032430, 2021.
\bibitem{he2016} K.\ He, X.\ Zhang, S.\ Ren y J.\ Sun. \emph{Deep residual learning for image
recognition}. CVPR, 2016.
\bibitem{bergholm2018} V.\ Bergholm et al. \emph{PennyLane: Automatic differentiation of hybrid
quantum-classical computations}. arXiv:1811.04968, 2018.
\end{thebibliography}

\end{document}
